% =====================================================================
%  Secao 1: Introdução
% =====================================================================
\section{Introdução}

Imagine que você está tentando sair de um labirinto.

Ao chegar a uma bifurcação, você escolhe um dos caminhos disponíveis. Caso encontre um beco sem saída, volta para a última bifurcação e tenta outra direção.

Essa estratégia parece natural para seres humanos. Curiosamente, ela também é utilizada por computadores para resolver diversos tipos de problemas.

A técnica responsável por esse comportamento é chamada de \emph{Backtracking}.

\begin{ideiachave}
O termo Backtracking pode ser traduzido livremente como ``voltar atrás''. Seu objetivo é explorar diferentes possibilidades até encontrar uma solução válida, retornando sempre que uma escolha realizada anteriormente impedir o progresso da busca.
\end{ideiachave}

\figura{fig:labirinto}%
  {Assim como uma pessoa em um labirinto, o algoritmo retorna ao último ponto de decisão quando encontra um caminho sem saída.}%
  {figuras/tikz/labirinto}
